当指尖落在键盘上，敲下这最后一章致谢时，心中百感交集。
既有完成毕业论文的释然与喜悦，也暗藏着即将与四年本科时光挥手告别的淡淡忧伤 —— 这是独属于毕业季的复杂心绪，难以言喻，却又无比真切。

在此，我要向我的导师致以诚挚的谢意。
您不仅是我学术道路上的引路人，更是我生活中的良师益友。
从最初的选题构思到最终的论文定稿，您全程悉心指导，带领我完成了人生中第一篇真正意义上的学术论文。
每当我在研究中遇到瓶颈与困惑，您总能以严谨的治学态度和耐心的讲解，与我一同剖析问题、寻找解决方案，引领我踏入科研的殿堂，拓宽了我的学术视野。
在生活中，您亦如兄长般关怀备至，球场上的并肩挥拍，餐桌上的欢声笑语，深夜里的促膝长谈，都让我在异乡求学的日子里感受到了别样的温暖。

感谢我的辅导员，从大一入学的第一次班会开始，您便一路陪伴着我们，走过了大学四年的春夏秋冬，直到我们站在毕业的十字路口，写下这篇致谢的终章。
是您耐心的引导，让我从初入校园时对大学生活的茫然无措，逐渐适应了独立的学习与生活节奏；是您细致的关怀，在我遇到学业困惑、生活难题时，总能第一时间为我指点迷津、排忧解难；是您真诚的鼓励，在我对未来感到迷茫时，帮我理清方向、重拾信心。
四年里，您不仅是我们的老师，更是我们的朋友和引路人，用责任与爱心，守护着我们从懵懂新生到成熟毕业生的每一步成长。

感谢我的班级，让我遇见了如此温暖又优秀的集体。
大一那年抱着试一试的心态加入，未曾想竟开启了一段长达四年的珍贵缘分。
我们曾在期末月并肩挑灯夜战，在假期的餐桌上分享欢声笑语，在升学择校的迷茫中彼此出谋划策，在开学选课的混乱里互通有无。
这个集体最动人的地方，在于它永远不是少数人的舞台，而是每个人都能找到自己位置的港湾。
我们相处融洽，不分彼此，谁遇到困难都会有人主动伸出援手；我们彼此成就，相互照亮，每个人的闪光点都能被看见、被珍视。
是你们身上那股不服输的拼劲赋予我勇往直前的动力，是你们不离不弃的陪伴帮助我撑过无数个难熬的时刻。

感谢我的舍友，感谢我们四年朝夕相伴的时光。
我们在生活中互帮互助、彼此包容，在琐碎的日常里积攒了最深厚的情谊。
疫情封校时的互相投喂，深夜组队开黑的酣畅淋漓，考试前互相划重点的默契，睡前卧谈会的幽默打趣，生病时端水送药的悉心照料…… 无数个细碎又温暖的瞬间，让这间不足二十平米的小屋，成为了我在求学时一个温暖安心的港湾。
在这里，不用伪装，不用拘谨，可以做最真实的自己。
是你们的陪伴与包容，让我学会了如何与人相处、如何分享、如何爱人。
这段同吃同住、同喜同悲的岁月，将是我一生都难以忘怀的珍贵回忆。

感谢我的女友，从大二与你相遇开始，我的大学生活便不再是孤身一人。
遇到困难，我们携手面对；陷入低谷，我们互相搀扶着慢慢走出；我们一起分享过彼此成功的喜悦，也共同分担过彼此失意的落寞。
但好在有你的陪伴与鼓励，让我在大学四年中慢慢成长，让我在求学路上始终有温暖的依靠，让我平淡的日子也变得闪闪发光。
你是我青春里最美好的遇见，也是我未来路上最坚定的同行者。

最后，我要将心底最深、最沉的谢意，献给我的父母与家人。
你们永远是我最坚实的后盾，最温暖的港湾。
二十余载含辛茹苦的养育，四年求学路上毫无保留的守护，你们的包容、无条件的支持与爱，是我能够心无旁骛完成学业、勇敢追逐梦想的根本动力。
是你们教会我热爱生活、善待世界、真诚待人；是你们在我迷茫困惑时为我点亮前行的灯塔，在我孤独无助时给予我最温暖的拥抱，在我疲惫憔悴时为我撑起一片可以安心休憩的天空。

回首望去，四年时光如白驹过隙。仿佛昨日我还是那个拖着行李箱踏入校园、对一切都充满新鲜与好奇的大一新生，转眼间我已站在毕业的十字路口，已然将要奔赴下一场山海。
这四年，说长不长，说短不短，却是我人生中最珍贵、最特殊的一段时光。太多的瞬间值得铭记，太多的回忆难以割舍，那些哭过笑过、拼过爱过的日子，早已化作生命中最温暖的印记。

还有太多的话想说，还有太多的事想做，但轻舟已过万重山，前路漫漫亦灿灿。
愿我们此去前程似锦，再相逢依旧如故。